% Interpretation of the tables produced by Rscript paperi/laske_taulukot.R.
\paragraph{Interpretation.}
The ordering in the tables matches the theorems. With the default settings,
\pol{pito} beat \pol{busiest} in 193 of 200 seeds, and the confidence interval
$[0.38,\ 0.44]$ of the paired difference excludes zero
(Theorem~\ref{la:aito}). In the settings $G = 1$ and $\eta = \varepsilon$ the
difference was exactly zero for every seed, as Theorem~\ref{la:valttamaton}
predicts.

Most of \pol{pito}'s advantage over \pol{rotation} comes from taking walking
volumes into account: \pol{busiest} removes about 2.1 expected falls and the
vibration detections about 0.4 more. The detections recover roughly a quarter
of the advantage of perfect information. With the default $\rho = 0$ the
slipperiness is redrawn at midnight, so the 00:00 round has no information about
the new day and the 06:00 round only a quiet night's worth. Mainly the 12:00 and
18:00 rounds are informed. When slipperiness persists from day to day
($\rho = 0.9$), information accumulates and the share rises to about 60 per cent.

Even the oracle removes only about 24 per cent of the falls, because the crew
covers 2 of the 20 blocks at a time. The setting $H = 12$ lies outside
Assumption~\ref{ol:A}. There \pol{pito} was still better on average, but the
theorems do not guarantee it (Remark~\ref{hu:H}).

The numbers are simulated falls in a $300 \times 300$ metre area over one month.
They are not insurance claims or real events in Turku.
